\documentclass[letterpaper,11pt,oneside,english]{scrreprt}

\usepackage{amsmath,amssymb}
\usepackage{geometry}
\usepackage{graphicx}
\usepackage[hidelinks]{hyperref}
\usepackage{cleveref}
\geometry{margin=1in}


\title{ADCS Concept of Operations (ConOps)}
\author{Sebastian Benner}
\date{\today}

\begin{document}
	
	\maketitle
	\tableofcontents
	\newpage
	
	%==================================================
	\chapter{Overview}
	
	The Attitude Determination and Control System (ADCS) is not continuously active. Instead, it operates on demand and is invoked either automatically or via commands from the ground.
	
	The system is designed to:
	
	\begin{itemize}
		\item Execute specific attitude control tasks when requested
		\item Transition between defined operational states
		\item Report completion and status back to the ground
	\end{itemize}
	
	All ADCS operations follow a structured pattern consisting of an initial state, a defined task, and a shutdown operation.
	
	%==================================================
	\chapter{Operational Philosophy}
	
	The ADCS operates as a task-based system rather than a continuously running control loop.
	
	\begin{itemize}
		\item The system is activated by either:
		\begin{itemize}
			\item Ground command
			\item Autonomous onboard trigger
		\end{itemize}
		
		\item Once activated, the ADCS performs a defined task
		
		\item Upon completion, the system:
		\begin{itemize}
			\item Sends a status message indicating completion
			\item Returns to an inactive or low-power mode if no further tasks are scheduled
		\end{itemize}
	\end{itemize}
	
	This approach minimizes power usage while maintaining full autonomy during task execution.
	
	%==================================================
	\chapter{Operational States}
	
	The ADCS operates within a set of discrete states.
	
	\section{Inactive State}
	
	\begin{itemize}
		\item ADCS control is not actively running
		\item Minimal power consumption
		\item Awaiting command from ground or autonomous trigger from C3
	\end{itemize}
	
	\section{Active Control State}
	
	\begin{itemize}
		\item ADCS estimation, guidance, and control loops are active
		\item Spacecraft is executing a defined task
	\end{itemize}
	
	\section{Completed State}
	
	\begin{itemize}
		\item Task objectives have been met
		\item Spacecraft has reached the desired attitude or condition
		\item Status message is transmitted to indicate completion
	\end{itemize}
	
		%==================================================
		
	\chapter{Operational Summary}
	
	All ADCS operations follow a consistent pattern:
	
	\begin{enumerate}
		\item Initial state (inactive or prior state)
		\item Task invocation (command or autonomous trigger)
		\item Task execution (control active)
		\item Task completion
		\item Status reporting
		\item Transition to end state (inactive or standby)
	\end{enumerate}
	
	This structure ensures predictable behavior, low power consumption, and clear communication of system status.
	
	\section{Nominal Operational Sequence}
	
	A typical operational sequence follows the order below:
	
	\begin{enumerate}
		\item Ground command is received specifying:
		\begin{itemize}
			\item Mission type
			\item Duration
			\item Whether data should be offloaded after completion
		\end{itemize}
		
		\item The ADCS is initialized:
		\begin{itemize}
			\item State estimation is initialized
			\item If necessary, the spacecraft performs a slow rotation about the +z axis until the star tracker is no longer occluded
		\end{itemize}
		
		\item The spacecraft executes the commanded mission for the specified duration
		
		\item Upon completion, the system evaluates upcoming ground station overpasses
		
		\item If the next overpass is sufficiently far in the future:
		\begin{itemize}
			\item The spacecraft enters a low-power mode
			\item ADCS control is shut down
		\end{itemize}
		
		\item Prior to the overpass window:
		\begin{itemize}
			\item The ADCS is reinitialized if previously shut down
			\item The spacecraft transitions to target tracking mode
		\end{itemize}
		
		\item Data is offloaded during the overpass
		
		\item After reaction wheel shutdown, a detumble operation is performed:
		\begin{itemize}
			\item Executed at a lower update rate
			\item Ensures the spacecraft is in an acceptable state for the next mission
			\item Results in a small residual spin about the +z axis
		\end{itemize}
		
		\item The system sends a final status message and transitions to the inactive state
	\end{enumerate}
	
	%==================================================
	
	\chapter{Detumble Operation}
	
	\section{Trigger Conditions}
	
	Detumble is initiated automatically following deployment.
	
	The trigger condition is:
	
	\begin{itemize}
		\item A predefined time delay after deployment
	\end{itemize}
	
	This ensures that the spacecraft begins stabilizing without requiring ground intervention.
	
	\section{Execution}
	
	\begin{itemize}
		\item Magnetorquer-based control is used to reduce angular velocity
		\item The system monitors angular rate until it falls below a defined threshold
	\end{itemize}
	
	\section{Completion}
	
	\begin{itemize}
		\item Once angular rates are sufficiently reduced:
		\begin{itemize}
			\item The detumble task is considered complete
			\item A status message is sent to the ground
		\end{itemize}
		
		\item The spacecraft transitions to a stable state and awaits further commands
	\end{itemize}
	
		%==================================================
	\chapter{Satellite Axis Definitions}
	
	OreSat1 and SENTINEL differ in their axis definitions.
	
	\section{OreSat1}
	
	\begin{figure}[h!]
		\centering
		\includegraphics[width=0.65\textwidth]{OreSat1_Axes_Deployed.png}
		\caption{OreSat1 axis definition}
		\label{fig:OreSat1}
	\end{figure}
	
	 As illustrated in \Cref{fig:OreSat1}, OreSat1 defines its axes as follows:
	
	\begin{itemize}
		\item +X axis points out along the same axis as the star tracker's boresight
		\item +Z axis points along the same axis along which the helical antenna is deployed, while -Z is the side of the satellite containing the CFC camera and UHF antennas
		\item +Y axis is derived from +X and +Z orientations using a standard orthogonal coordinate system
	\end{itemize}
	
	\section{SENTINEL}
	
	No image similar to that seen in \Cref{fig:OreSat1} is available to illustrate the SENTINEL axis definitions at the time of writing, which are similar to those of OreSat1, but rotated about the +Y axis by 180 degrees:
	
	\begin{itemize}
		\item +X axis points along the same axis (face???) as the Radio Occultation (RO) antenna, while the -X axis points out along the same axis as the star tracker's boresight
		\item +Z axis points out of the side of the satellite on which the UHF antenna are mounted, while the -Z points along the same axis along which the helical antenna is deployed
		\item +Y axis is derived from +X and +Z orientations using a standard orthogonal coordinate system
	\end{itemize}
	
	%==================================================
	\chapter{Mission Execution}
	
	\section{Commanded Operation}
	
	Mission execution is initiated via ground command.
	
	The command specifies:
	
	\begin{itemize}
		\item Mission type (e.g. nadir pointing, ram-facing orientation)
		\item Duration of execution
		\item Whether data should be downlinked after completion
	\end{itemize}
	
	\section{Execution Behavior}
	
	Once commanded:
	
	\begin{itemize}
		\item The ADCS transitions from the inactive state to the active control state
		\item The spacecraft performs the requested mission for the specified duration
	\end{itemize}
	
	\section{Attitude Modes}
	
	\begin{itemize}
		\item \textbf{Nadir Pointing:} Spacecraft +z axis aligned with Earth center
		\item \textbf{Ram-Facing:} Spacecraft +x axis aligned with velocity vector
	\end{itemize}
	
	For circular or near-circular orbits, these modes are nearly equivalent. In elliptical orbits, the difference becomes more pronounced due to variation in the velocity vector.
	
	\section{Completion}
	
	\begin{itemize}
		\item After the specified duration:
		\begin{itemize}
			\item The task is terminated
			\item A status message is transmitted
		\end{itemize}
		
		\item The spacecraft transitions to the completed or inactive state
	\end{itemize}
	
	%==================================================
	\chapter{Autonomous Downlink Operation}
	
	\section{Overview}
	
	If data downlink is requested, the spacecraft autonomously schedules communication following mission completion.
	
	\section{Overpass Prediction}
	
	\begin{itemize}
		\item The spacecraft predicts upcoming ground station overpasses
		\item The next suitable overpass window is selected
	\end{itemize}
	
	\section{Pre-Overpass Behavior}
	
	\begin{itemize}
		\item The spacecraft remains in a low-power or standby state
		\item ADCS is not actively performing high-rate control during this period
	\end{itemize}
	
	\section{Communication Execution}
	
	\begin{itemize}
		\item Shortly before the overpass window:
		\begin{itemize}
			\item ADCS is activated
			\item The spacecraft transitions to target tracking mode
		\end{itemize}
		
		\item The spacecraft points toward the ground station
		
		\item Data transmission begins once the ground station is within range
	\end{itemize}
	
	\section{Completion}
	
	\begin{itemize}
		\item After the overpass:
		\begin{itemize}
			\item Transmission stops
			\item A status message is generated
		\end{itemize}
		
		\item The spacecraft returns to an inactive or standby state
	\end{itemize}
	
	
\end{document}